% TOP_PROBLEM: {"id": 3031, "title": "Graphs of exact colorings", "category_id": 4, "category": {"id": 4, "name": "algebra", "display_name": "Algebra", "description": "Group theory, ring theory, field theory, and algebraic structures.", "slug": "algebra", "order_index": 4, "created_at": "2026-07-31T15:26:25.671Z"}, "status": "open", "problem_number": "OPG-57824", "set_id": 12}
% FIRST_POSTED: 2026-10-01
% ATTEMPT_STATUS: unresolved
% UnsolvedMath ID 3031; source catalogue code OPG-57824
% !TeX program = lualatex
% Research attempt by GPT-6 Astra Ultra; no claim of a new complete solution.
\documentclass[11pt,a4paper]{article}
\usepackage[margin=25mm]{geometry}
\usepackage{fontspec}
\setmainfont{lmroman10-regular.otf}[BoldFont=lmroman10-bold.otf,
 ItalicFont=lmroman10-italic.otf,BoldItalicFont=lmroman10-bolditalic.otf]
\usepackage{amsmath,amssymb,mathtools}
\usepackage{unicode-math}
\setmathfont{latinmodern-math.otf}
\AtBeginDocument{\let\setminus\smallsetminus}
\usepackage{xurl,hyperref}
\hypersetup{colorlinks=true,urlcolor=blue}
\setcounter{secnumdepth}{0}
\setlength{\parindent}{0pt}
\setlength{\parskip}{5pt}
\setlength{\emergencystretch}{3em}
\begin{document}
\section*{Graphs of exact colorings}\label{3031}
{\small UnsolvedMath ID 3031 · OPG-57824 · Algebra · OpenGarden}
\subsection{Definitions and mathematical statement}
Let $\mathbb N=\{1,2,\ldots\}$ and let $[\mathbb N]^2$ be its set of unordered
two-element subsets, the edge set of the countably infinite complete graph.
For integers $c\ge m\ge1$, an exact $c$-colouring is a surjection
$\varphi:[\mathbb N]^2\to\{1,\ldots,c\}$. There is no proper-colouring
condition: edges sharing a vertex may have the same colour.

For an infinite subset $X\subseteq\mathbb N$, write
$\mathcal C_\varphi(X)=\{\varphi(e):e\in[X]^2\}$.
The complete subgraph on $X$ is exactly $m$-coloured when
$|\mathcal C_\varphi(X)|=m$. Define $P(c,m)$ to mean that every exact
$c$-colouring admits an infinite $X$ with this property.
The problem asks whether the following classification holds for all the
specified integer pairs:
\[
 P(c,m)\quad\Longleftrightarrow\quad
 m=1\ \text{or}\ m=2\ \text{or}\ m=c.
\]
The positive assertion is universal over colourings, whereas a negative
answer for a particular pair requires only one surjective colouring for
which every infinite vertex set uses a number of colours different from $m$.
Colours must all occur globally, but no condition requires a colour to occur
infinitely often or at infinitely many vertices. The desired subgraph must
have infinitely many vertices; finding finite complete subgraphs with the
required count does not suffice. The problem concerns exact colour counts,
so the stronger demand that a particular prescribed collection of $m$
colours occur is not part of the statement.
\subsection{Short English statement}
Which exact colour counts must occur on an infinite complete subgraph of every finite edge-colouring of the countable complete graph?
\subsection{Research attempt}
\paragraph{Scope and process.}
This is a GPT-6 Astra Ultra research attempt dated 1 October 2026, using
primary-source browsing and elementary mathematical checks. The canonical
definition above is retained. Results below are restricted results or
reductions, not a claim of a new solution to the entire catalogue question.
No formal proof assistant or external mathematical referee checked this text.


\paragraph{Literature check.}
The exact-colour-count conjecture is attributed to Erickson in [S2].
Narayanan [R1] studies the same infinite-subgraph spectrum and explains
why guaranteed counts and the number of available counts are different
questions. The arguments below check the positive cases directly and
give an elementary family of excluded intermediate counts. These are
known benchmark constructions, not a new classification.

\paragraph{The three forced cases.}
The case $m=c$ uses $X=\mathbb N$. For $m=1$, the infinite Ramsey
argument can be made explicit: choose $v_1$, then an infinite set on
which all its outgoing edges have one colour; choose $v_2$ there and
repeat within a nested infinite set. Infinitely many of the selected
vertices have the same outgoing colour, by the infinite pigeonhole
principle. Their induced complete graph is monochromatic.

Now suppose $c\ge2$. Choose such an infinite monochromatic set $Y$,
of colour $a$. Some edge $uv$ has colour $b\ne a$, by surjectivity.
Delete $u,v$ from $Y$ if necessary. Pass to an infinite subset $Z$ of
$Y$ on which the ordered pair of colours on $uz,vz$ is constant,
say $(p,q)$; only finitely many pairs are possible. If $p\ne a$, the
set $Z\cup\{u\}$ uses precisely $a,p$. If $q\ne a$, use $v$
instead. If both equal $a$, the set $Z\cup\{u,v\}$ uses precisely
$a,b$. Hence $P(c,2)$ holds whenever it is in the allowed range.
No colour class is assumed infinite: a single exceptional edge suffices.

\paragraph{A counterexample mechanism for infinitely many pairs.}
Take $h\ge3$, put $c=1+\binom h2$, and designate a finite set $B$
of $h$ vertices. Give each edge inside $B$ a different private colour,
and give all other edges one common colour $0$. Relabel the colours
by $[c]$ if desired. For any infinite $X$, let $t=|X\cap B|$.
The common colour must occur, since $X\setminus B$ is infinite,
and exactly $\binom t2$ private colours occur. Conversely every
$0\le t\le h$ is attained by choosing that many vertices of $B$
and infinitely many outside. Thus the spectrum is exactly
\[
 \left\{1+\binom t2:0\le t\le h\right\}.
\]
For $c=4$ it is $\{1,2,4\}$, so $P(4,3)$ is false. For $c=7$
it is $\{1,2,4,7\}$, excluding $m=3,5,6$. In general every
intermediate $m$ outside the displayed set is excluded for this $c$.

\paragraph{Attempted interpolation and its failure.}
One might add independent private colours to increase $c$ while preserving
a missing count. This need not work. For example, begin with the
four-colour example on a private triangle. Add one new private edge on
two new vertices disjoint from that triangle. The resulting exact
five-colouring has a three-coloured infinite subgraph: take two triangle
vertices, both new vertices, and infinitely many ordinary vertices.
It sees one triangle colour, the new colour and colour $0$. Thus even
this minimal padding destroys the exclusion of $m=3$. A classification
for arbitrary $(c,m)$ needs constructions whose spectra remain controlled
under the required change of total colour count.

\paragraph{Verification.}
The Ramsey extraction only uses countably many nested infinite sets;
no assertion that their intersection is infinite is made. The chosen
subsequence, rather than the intersection, is the monochromatic set.
In the two-colour argument both special vertices may originally lie
outside $Y$ or just one may do so; deletion and thinning handle both
cases. The finite exceptional set construction concerns infinite
subgraphs, so its unavoidable background colour cannot be discarded
as it could for a finite subgraph wholly inside $B$.

\paragraph{Final status and first remaining gap.}
\textbf{UNRESOLVED.} The universally positive cases are verified, as
are infinitely many negative parameter pairs. There is no construction
here for every $3\le m<c$; the interpolation obstruction identifies
the first unsupported step in trying to turn this family into a full proof.

\subsection{Sources}
\begin{itemize}
\item[S1] ULAM.AI and the UnsolvedMath Contributors, supplied problems.json, record 3031 (OPG-57824), OpenGarden collection. Catalogue identity and source transcription; CC BY 4.0.
\url{https://huggingface.co/datasets/ulamai/UnsolvedMath}.
\item[S2] Open Problem Garden, Graphs of exact colorings. Original problem and discussion; the mathematical formulation below is expanded from the supplied source record.
\url{http://www.openproblemgarden.org/op/graphs_of_exact_colorings}.

\item[R1] Bhargav Narayanan, \emph{Exactly $m$-coloured complete infinite
subgraphs}, research paper, 2013 preprint. Sections 1--2 checked for the
formulation and small-rainbow construction.
\url{https://sites.math.rutgers.edu/~narayanan/pdf/m_col.pdf}.

\end{itemize}
\end{document}
